% ===== 07-discussion.tex =====
\section{Discussion and Limitations}
\label{sec:discussion}
\textbf{Asynchronous tools.} Thread-keyed context fails silently for coroutines that yield between markers (E4), and no kernel signal can detect it inside one thread. Binding the identity to the task and handing it to the kernel at the system call that creates the connection or process removes the failure (E5) without changing tool code. The approach relies on the runtime propagating task context, which Python's \texttt{asyncio} and \texttt{anyio} provide through context variables (\texttt{contextvars}); equivalent mechanisms exist in other runtimes (e.g., asynchronous local storage in Node.js), but we implemented and evaluated only Python.

\textbf{Threat model and trust.} The identity is asserted by the agent runtime, which we trust as the operator's component; tools and MCP clients are not trusted. In the prototype, this trust is not enforced: any process can emit a marker, and the \texttt{\_meta} identity is chosen by the client, so a malicious tool or client could claim another execution's identity. The policy is therefore a containment control for correct or buggy runtimes, not a security boundary against adversarial tools. Hardening is straightforward to specify but not evaluated here: accept markers only from registered runtime threads, freeze a process's identity at \texttt{exec}, and derive the identity from an authenticated MCP session instead of \texttt{\_meta}. The fused resolver bounds a faulty label only where the process structure is informative, and the policy's lineage check enforces the same bound in the kernel; in shared workers the remaining defenses are outside the kernel, e.g., auditing the runtime or isolating executions in separate processes. Deciding by the tree owner assumes that a registered tree is never handed over to another execution; denying on conflict makes no such assumption, which is why it is the default. A related limitation is that the per-thread entry is not cleared after a call, so a connection opened by native code that bypasses the wrappers inherits the last identity set on that thread.

\textbf{Threats to validity.} We follow the categories of Wohlin et al.~\cite{wohlin2012experimentation}. \emph{Internal:} capture loss, socket-pointer reuse in the connect/state join, and PID reuse (not exercised). \emph{Construct:} synthetic executions approximate, but are not, tasks or tool calls; attribution is evaluated for TCP connection attempts, not bytes, and pooled or reused connections (e.g., HTTP keep-alive shared across sessions) are not exercised; such reuse would credit later requests to the session that opened the connection. \emph{External:} a single kernel and host, a VMware virtual machine without CPU frequency control; loopback traffic; IPv4 only; one MCP SDK version; a single small local model in E7, whose tool-call behavior is part of the workload rather than a controlled factor. \emph{Conclusion:} events within a run are dependent, so statistics are computed over runs.

\textbf{Scope.} We attribute and police TCP connection attempts; byte-level accounting per execution and policies over destinations other than ports are natural extensions of the same identity. Like the accounting, the policy relies on the runtime-provided identity only where lineage cannot contradict it, i.e., inside a shared worker; there, attribution is runtime-asserted rather than kernel-verified and should be reported as such.
